%% kapa: exam=1 round=28 year=2017 period=2 form=A subject=appraisal_law subject_ko=감정평가관계법규 seq=22 number=22
%% source: Q-Net 감정평가사 기출문제 (cst00309) artlSeq=5199603 file=extracted/1008/5199603/2017년 제28회 감정평가사 1차시험 문제지/2교시 A형.pdf
%% conversion: kordoc 4.24.51 -> tools/gen1.py (auto); review=proofread-sonnet; flags=none
\begin{question}{22}
국유재산법령상 일반재산의 처분가격에 관한 설명으로 옳은 것은?
\begin{choices}
  \item 지식재산을 처분할 때의 예정가격은 두 개의 감정평가업자의 평가액을 산술평균한 금액으로 한다.
  \item 상장법인이 발행한 주권을 처분할 때의 예정가격은 하나의 감정평가업자의 평가액으로 한다.
  \item 비상장법인이 발행한 지분증권을 처분할 때의 예정가격은 두 개의 감정평가업자의 평가액을 산술평균한 금액으로 한다.
  \item 대장가격이 3천만원인 부동산을 공공기관에 처분할 때의 예정가격은 하나의 감정평가업자의 평가액으로 결정한다.
  \item 증권을 제외한 일반재산을 처분할 때의 예정가격에 대한 감정평가업자의 평가액은 평가일부터 3년 이내에만 적용할 수 있다.
\end{choices}
\end{question}
\answer{4}
